%%=============================================================================
%% Methodologie
%%=============================================================================

\chapter{\IfLanguageName{dutch}{Methodologie}{Methodology}}%
\label{ch:methodologie}

%% TODO: In dit hoofstuk geef je een korte toelichting over hoe je te werk bent
%% gegaan. Verdeel je onderzoek in grote fasen, en licht in elke fase toe wat
%% de doelstelling was, welke deliverables daar uit gekomen zijn, en welke
%% onderzoeksmethoden je daarbij toegepast hebt. Verantwoord waarom je
%% op deze manier te werk gegaan bent.
%% 
%% Voorbeelden van zulke fasen zijn: literatuurstudie, opstellen van een
%% requirements-analyse, opstellen long-list (bij vergelijkende studie),
%% selectie van geschikte tools (bij vergelijkende studie, "short-list"),
%% opzetten testopstelling/PoC, uitvoeren testen en verzamelen
%% van resultaten, analyse van resultaten, ...
%%
%% !!!!! LET OP !!!!!
%%
%% Het is uitdrukkelijk NIET de bedoeling dat je het grootste deel van de corpus
%% van je bachelorproef in dit hoofstuk verwerkt! Dit hoofdstuk is eerder een
%% kort overzicht van je plan van aanpak.
%%
%% Maak voor elke fase (behalve het literatuuronderzoek) een NIEUW HOOFDSTUK aan
%% en geef het een gepaste titel.


\subsection {Fase 1: Iteratie 1 - De klassieke dreiging (LKM & Singularity)}
Het doel van de eerste iteratie is het selecteren van en het verder verdiepen in een klassieke Loadable Kernel Module (LKM) rootkit. Deze zal op een educatieve manier ontmanteld worden en op een gestructureerde manier uitgelegd worden aan de hand van de literatuurstudie die tijdens de eerste iteratie opgesteld wordt. De onderliggende methodes en technieken moeten eerst grondig begrepen worden vooraleer ze op een simpele manier kunnen worden uitgelegd. Daarom is het belangrijk om een eenvoudige rootkit te zoeken die wel nog genoeg functionaliteiten heeft heeft om alle methodes en technieken toe te kunnen lichten. Als de gevonden rootkit volstaat aan de bovenstaande criteria zal er een virtuele omgeving opgezet worden als een proof of concept. Dit is deels voor de veiligheid maar ook voor het gebruiksgemak. Het opzetten van de omgeving geeft de mogelijkheid om de functionaliteiten van de rootkit te testen. Als de rootkit een Command & Control server functionaliteit bevat zou dit nuttig kunnen zijn vanuit een educatief oogpunt. Verder moeten de hooking technieken begrepen worden en hoe deze werken op het systeem. De ondervonden technieken moeten uitgewerkt worden in de literatuurstudie. Tijdens dit proces kunnen er ook andere methodes aan bod komen die niet direct gebruikt worden door de geselecteerde rootkit maar toch deel uitmaken van het moderne landschap van rootkits. Als alle concepten begrepen zijn moet er gekeken worden naar hoe dit zo goed mogelijk op een educatieve manier uitgelegd kan worden. Het eindresultaat van deze fase is een .ova-bestand van de gekozen LKM rootkit en kennis over hoe het werkt. Door de testopstelling op deze manier op te slaan, wordt gegarandeerd dat het experiment voor de demonstratie volledig reproduceerbaar en repliceerbaar is.

\subsection {Fase 2: Iteratie 2 - De moderne dreiging (eBPF & bad-bpf)}
Uit de literatuurstudie van de vorige iteratie zal naar verwachting blijken dat het landschap van rootkits evolueert en er een nieuwe, moderne variant actief wordt gebruikt. De kernel wordt continu bijgewerkt door de ontwikkelaars. Dit zorgt ervoor dat er nieuwe methodes ontwikkeld worden, die bedoeld zijn om bedrijven en gebruikers nieuwe mogelijkheden te geven. Maar een negatief aspect hiervan is dat deze ook door rootkits gebruikt kunnen worden. In de tweede iteratie moet gezocht worden naar een eBPF rootkit (zoals bad-bpf) om te analyseren en demonstreren. Een nieuwe methode betekent ook dat nieuwe concepten aan bod komen die verder uitgewerkt moeten worden in de literatuurstudie. Hoewel eBPF-rootkits gebruik maken van zelfde basis concepten als een LKM zijn er toch verschillen. Om deze moderne bedreiging te testen zal de rootkit opnieuw geïsoleerd worden in een virtuele omgeving. Hierin kan er weer gekeken worden naar de functionaliteiten van deze rootkit, en de basisconcepten ervan moeten op een educatieve manier worden uitgelegd. De gevonden rootkit moet ook in een .ova-bestand gezet worden. 

\subsection {Fase 3: Synthese en Preventiestrategieën}
 Om het onderzoek te vervolledigen, moet ook getoond worden hoe je je tegen deze rootkits kunt verdedigen. Er zal via een literatuurstudie gezocht worden naar manieren om deze rootkits te detecteren en preventie maatregelen te nemen. De twee virtuele omgevingen die gemaakt zijn in de vorige iteraties zullen gebruikt worden om de technieken die verzameld zijn tijdens de literatuurstudie te testen. Als deze tests geslaagd zijn, moet dit tenslotte gedocumenteerd worden zodat het kan gebruikt worden op een educatieve manier. Het eindresultaat van deze iteratie zijn methodes om rootkits te detecteren en het besmetten van een kernel te voorkomen. Om dit eenvoudig te presenteren zal er een stappenplan gemaakt moet worden over hoe je eBPF en LKM rootkits kunt detecteren en voorkomen.